\corrigesChapitre{6}{Primitives et équations différentielles}
\begin{corrige}{1}{Primitives usuelles}
\begin{questions}
\item $x^3-4x+C$, $C\in\R$.
\item $\tfrac12\mathrm e^{2x}+C$.
\item $\ln x+C$. Chaque réponse se contrôle par dérivation.
\end{questions}
\end{corrige}
\begin{corrige}{2}{Reconnaître une composée}
\begin{questions}
\item $\mathrm e^{x^2}$.
\item $\ln(x^2+3)$, puisque le dénominateur est toujours positif.
\item $(x^2+1)^2$, dont la dérivée est $4x(x^2+1)$.
\end{questions}
\end{corrige}
\begin{corrige}{3}{Condition initiale}
\begin{questions}
\item $F=x^2+C$ ; $C=3$, donc $F=x^2+3$.
\item $G=\mathrm e^x+C$ ; $C=-\mathrm e$, donc $G=\mathrm e^x-\mathrm e$.
\item Une même dérivée correspond à toute une famille. La valeur imposée sélectionne un membre précis.
\end{questions}
\end{corrige}
\begin{corrige}{4}{Vérifier une solution}
\begin{questions}
\item $y\,\!^{\prime}=-6\mathrm e^{-2t}$.
\item $-2y+8=-8-6\mathrm e^{-2t}+8=-6\mathrm e^{-2t}$, égal à la dérivée.
\item $y(0)=7$.
\end{questions}
\end{corrige}
\begin{corrige}{5}{Résoudre sans terme constant}
\begin{questions}
\item $y(t)=2\mathrm e^{3t}$.
\item $z(t)=10\mathrm e^{-t/2}$.
\item La première tend vers $+\infty$ ; la seconde vers $0$, car les signes des coefficients exponentiels diffèrent.
\end{questions}
\end{corrige}
\begin{corrige}{6}{Déplacer un équilibre continu}
\begin{questions}
\item $\ell=12$.
\item $y(t)=12-10\mathrm e^{-t/2}$.
\item $\mathrm e^{-t/2}<0{,}1$ donne $t>2\ln10\simeq4{,}605$ unités de temps.
\end{questions}
\end{corrige}
\begin{corrige}{7}{Cas $a=0$}
\begin{questions}
\item Elle diviserait par zéro ; l’équation devient simplement $y\,\!^{\prime}=5$.
\item $y(t)=5t+C$ ; $9=10+C$, donc $y(t)=5t-1$.
\item $y(2)=10-1=9$ et la dérivée est bien $5$.
\end{questions}
\end{corrige}
\begin{corrige}{8}{Euler à la main}
\begin{questions}
\item $y_{n+1}=y_n-0{,}5y_n=0{,}5y_n$.
\item $y_1=0{,}5$, $y_2=0{,}25$.
\item La valeur exacte est $\mathrm e^{-1}\simeq0{,}367879$ ; Euler sous-estime d’environ $0{,}117879$.
\end{questions}
\end{corrige}
\begin{corrige}{9}{Pourquoi une seule constante ?}
\begin{questions}
\item $(F-G)^{\prime}=0$.
\item Une dérivée nulle sur un intervalle donne une constante : $G=F+C$.
\item Sur $\R\setminus\{0\}$, la fonction égale à $0$ pour $x<0$ et à $1$ pour $x>0$ a une dérivée nulle partout dans son domaine sans être une constante unique.
\end{questions}
\end{corrige}
\begin{corrige}{10}{Pas trop grand}
\begin{questions}
\item Le facteur vaut $-1$ ; les termes alternent entre $1$ et $-1$.
\item La solution exacte $\mathrm e^{-2t}$ reste positive et tend vers $0$. Euler avec ce pas donne un comportement artificiel.
\item Le facteur vaut $0{,}5$ ; les termes restent positifs et tendent vers zéro, en cohérence qualitative avec la solution exacte.
\end{questions}
\end{corrige}
\begin{corrige}{11}{Logistique : vérification guidée}
\begin{questions}
\item $y(0)=20$ et $y(t)\to100$.
\item $y\,\!^{\prime}=400\mathrm e^{-t}/(1+4\mathrm e^{-t})^2$ ; remplacer $y$ dans le second membre donne la même expression.
\item On contrôle une fonction déjà fournie. On n’a pas établi comment trouver toutes les solutions de cette équation non linéaire.
\end{questions}
\end{corrige}
\begin{corrige}{12}{Modèles discret et continu}
\begin{questions}
\item $\mathrm e^{-k}=0{,}8$, donc $k=-\ln(0{,}8)$.
\item $y\,\!^{\prime}=-k(y-50)$.
\item Non : Euler donnerait le coefficient $1-k=1+\ln(0{,}8)\simeq0{,}776856$, différent de $0{,}8$. Coïncidence aux temps entiers et approximation d’Euler sont deux constructions distinctes.
\end{questions}
\end{corrige}
